\begin{abstract}
Under a fixed task-training budget, does output-channel rotation affect
adaptation differently across parameterizations? We also test whether a shared
trainable orthogonal output map reduces Factor's Native--Rotated score gap.
We compose adapter residuals through frozen protein predictors' native
pair-generating operators and train native and rotated heads separately under
paired settings, starting from the same zero-residual function. Orthogonal controls preserve
local writer spectra; a generic head with explicit factor access and a free
output layer provides a function class closed under output rotation.
Under fixed 1536-update recipes, four-cell studies find greater rotation
sensitivity for Factor than G+ in specified OpenFold and Protenix settings.
On 192 prospectively locked new targets, fixed Protenix/ESMC models establish
a C$\alpha$ pair-lDDT interaction of \result{p192_pair_interaction}
[\result{p192_pair_interactionLo}, \result{p192_pair_interactionHi}], whereas
OpenFold's Fresh96 primary interaction remains unestablished.
Within Factor, a shared orthogonal compensator improves two preselected
OpenFold rotated arms. Its extra benefit over Native is
\result{e2_pair_T_C} in the original execution but
\result{e2rep_pair_T_C} in a same-seed retraining repeat: rotated-arm benefit
is observed in both, but relative mitigation was not re-established in the repeat.
A prespecified reconstruction diagnostic did not establish the predicted ranking
of post-training rotation penalties across eight tested rotations.
The new-target test supports parameterization-dependent rotation sensitivity
for one fixed recipe; useful output freedom does not ensure a repeated extra
benefit over Native. Adaptation gains, rotation sensitivity and head
competitiveness therefore require separate evaluation.
\end{abstract}
